%%%%%%%%%%%%%%%%%%%%%%%%%%%%%%%%%%%%%%%%%%%%%%%%%%%%%%%%%%%%%%%%%%%%%%%%%%%%
% Canadian Climate Framing (CCF) -- Supplementary Information
%
% This document is the standalone Supplementary Information that
% accompanies CCF_Methodology.tex. Tables are numbered with the
% prefix ``S'' (S1, S2, ..., S12).
%
% This LaTeX source is maintained by hand. Tables S1, S2, S3 and S6
% are inlined verbatim below (their content is prose / static lists
% that do not depend on the live data). Tables S4 to S12 are
% reproducibly generated by ``Scripts/Annotation/17_generate_tables.py``
% from the canonical CSVs in ``Database/Training_data/`` and the
% live ``CCF_Database`` (S8 only); each table is written to
% ``paper/CCF_Methodology/Results/Outputs/Tables/`` and \input{}ed
% from the appropriate landscape block below.
%
% Cross-reference convention. The main manuscript refers to tables
% in this file as ``Supplementary Table SX'' using literal numbers;
% the numbering is fixed by the order of the blocks below.
%%%%%%%%%%%%%%%%%%%%%%%%%%%%%%%%%%%%%%%%%%%%%%%%%%%%%%%%%%%%%%%%%%%%%%%%%%%%

\documentclass[12pt]{article}

\usepackage[utf8]{inputenc}
\usepackage[T1]{fontenc}
\usepackage[english]{babel}
\usepackage{geometry}
\geometry{margin=2.5cm}
\usepackage{booktabs}
\usepackage{longtable}
\usepackage{ltcaption}
\usepackage{multirow}
\usepackage{array}
\usepackage{tabularx}
\usepackage{colortbl}
\usepackage{graphicx}
\usepackage[table]{xcolor}
\usepackage{authblk}
\usepackage{pdflscape}
\usepackage{afterpage}
\usepackage{float}

% pdflscape rotates the page but leaves \textwidth at its portrait
% value (16.6 cm on Letter with 2.5 cm margins), which causes
% longtables to be left-aligned against the portrait text block
% rather than centred on the rotated page. The wrapper below grows
% \textwidth / \linewidth / \hsize to the rotated page width
% (\paperheight - 5 cm = 22.0 cm) for the duration of the landscape
% block, and longtable's default centring then places every table
% in the middle of the rotated page.
\makeatletter
\let\CCFlandscape@orig\landscape
\renewcommand{\landscape}{%
  \CCFlandscape@orig
  \setlength{\textwidth}{\dimexpr\paperheight-5cm\relax}%
  \setlength{\linewidth}{\textwidth}%
  \setlength{\hsize}{\textwidth}%
  \setlength{\columnwidth}{\textwidth}%
}
\makeatother
% Centre longtables horizontally inside \linewidth (works in concert
% with the wrapper above, which sets \linewidth to the rotated width
% during landscape blocks).
\AtBeginDocument{%
  \setlength{\LTleft}{0pt plus 1fill}%
  \setlength{\LTright}{0pt plus 1fill}%
}

% Default caption width for longtables.
\setlength{\LTcapwidth}{\linewidth}

\usepackage{csquotes}
\usepackage[backend=biber,style=numeric,sorting=none]{biblatex}
\addbibresource{references.bib}

% xurl must be loaded BEFORE hyperref so it can patch \url's tokeniser
% to allow line breaks at every character of a long URL. Without it,
% Zenodo, HuggingFace, GitHub, and Wikipedia links in the
% bibliography routinely overflow the right margin by 60-90pt.
\usepackage{xurl}
\usepackage[
  unicode=true,
  colorlinks=true,
  linkcolor=blue,
  citecolor=blue,
  urlcolor=blue,
  breaklinks=true
]{hyperref}
\usepackage{cleveref}

% \sloppy globally in the bibliography so biblatex's URL/DOI/note
% fields do not push lines past the right margin when xurl cannot
% find a natural break point.
\appto\bibsetup{\sloppy}

% Relax the tolerance for the body text so long paragraphs containing
% \texttt{...} identifiers and bilingual quotes do not exhaust TeX's
% default break-finding budget. Same values as in the main manuscript.
\tolerance=800
\emergencystretch=3em

\crefname{table}{Supplementary Table}{Supplementary Tables}
\Crefname{table}{Supplementary Table}{Supplementary Tables}

\renewcommand\_{\texttt{\detokenize{_}\allowbreak}}

% All tables in this document use the ``S'' prefix.
\renewcommand{\thetable}{S\arabic{table}}
\renewcommand\theHtable{Supplementary.\thetable}

\setcounter{secnumdepth}{-1}  % no section numbering (Scientific Data requirement)

% ---------------------------------------------------------------------------
% Corpus-level figures — SINGLE SOURCE OF TRUTH for the manuscript AND the SI.
% Update these macros (and regenerate the input tables/figures) when a new
% database version is released; do not hard-code corpus numbers in the text.
% ---------------------------------------------------------------------------
\newcommand{\CCFarticles}{283{,}964}      % articles in the corpus
\newcommand{\CCFsentences}{9{,}908{,}776} % two-sentence analytical units
\newcommand{\CCFsentApprox}{9.9}          % ditto, millions (approx, one decimal)
\newcommand{\CCFembeddings}{10{,}192{,}740}% rows of CCF_sentence_embeddings
\newcommand{\CCFembedApprox}{10.19}        % ditto, millions (approx)
\newcommand{\CCFnewspapers}{22}           % distinct newspaper titles
\newcommand{\CCFyearfirst}{1978}
\newcommand{\CCFyearlast}{2026}
\newcommand{\CCFnyears}{48}               % elapsed span: yearlast - yearfirst
                                          % (Feb 1978 -> Jul 2026), not the count
                                          % of calendar years touched
\newcommand{\CCFenshare}{83.4}            % share of English units (%)
\newcommand{\CCFfrshare}{16.6}            % share of French units (%)
% Version-specific deposit identifiers and sizes (update on each release)
% One version number for all three deposits (data SQL, data Parquet, code):
% they are released together and must never drift apart again.
\newcommand{\CCFversion}{2.1.0}
\newcommand{\CCFsqldoi}{10.5281/zenodo.22019287}     % PostgreSQL edition (version DOI)
\newcommand{\CCFparquetdoi}{10.5281/zenodo.22019288} % Parquet mirror (version DOI)
\newcommand{\CCFcodedoi}{10.5281/zenodo.22019289}     % code & documentation deposit
\newcommand{\CCFsqldoiConcept}{10.5281/zenodo.20346363}
\newcommand{\CCFparquetdoiConcept}{10.5281/zenodo.20346372}
\newcommand{\CCFcodedoiConcept}{10.5281/zenodo.21921214}
\newcommand{\CCFsqlsize}{40}              % GB, CCF_Database.tar
\newcommand{\CCFparquetsize}{17.5}        % GB, Parquet mirror

\title{\textbf{CCF Database: A Machine-Learning-Annotated Corpus of \CCFarticles{} Canadian Climate Articles (\CCFyearfirst{}--\CCFyearlast{})}\\[0.4em]
\large Supplementary Information}

\author[1]{Antoine Lemor\thanks{antoine.lemor@umontreal.ca}}
\author[2]{Aliz\'ee Pillod\thanks{alizee.pillod@umontreal.ca}}
\author[2]{Matthew Taylor\thanks{matthew.taylor@umontreal.ca}}
\author[2]{Richard Nadeau\thanks{richard.nadeau@umontreal.ca}}
\affil[1]{Universit\'e de Sherbrooke}
\affil[2]{Universit\'e de Montr\'eal}

\date{}

\begin{document}
\maketitle
\thispagestyle{empty}

\noindent This document collects the Supplementary Information accompanying the CCF Database manuscript. It contains: prior-literature and corpus composition tables (Supplementary Tables~S1 and~S2), the complete annotation framework with operational definitions (Supplementary Table~S3), the complete per-category training metrics and dataset distribution (Supplementary Tables~S4--S5), the Named Entity Recognition performance summary (Supplementary Table~S6), the per-category validation metrics on the 1{,}000-sentence gold standard (Supplementary Table~S7), the per-category distribution of annotations across the database (Supplementary Table~S8), and four additional tables added during revision: per-category inter-coder reliability (Supplementary Table~S9), reliability tier assignment (Supplementary Table~S10), per-model training provenance (Supplementary Table~S11), and the complete data dictionary of the enriched database (Supplementary Table~S12).

\vspace{0.5em}
{\sloppy
\noindent\textit{Reproducibility note.} Tables~S1, S2, S3, and~S6 are static prose / lookup tables inlined directly in this LaTeX source (their content does not depend on the live data). Tables~S4 to~S12 are reproducibly generated by \texttt{Scripts/Annotation/17\_generate\_tables.py} from the canonical CSVs in the \texttt{Database/Training\_data/} directory and the live \texttt{CCF\_Database} (Table~S8 only); the generator reuses the category-normalisation API exposed by \texttt{Scripts/Annotation/15\_normalize\_categories.py}. Re-running scripts~15, 16, and~17 followed by a fresh \texttt{pdflatex} pass regenerates every numerical value in this document from the underlying data.
\par}

%==========================================================================
% Supplementary Table S1 -- Previous studies on climate change coverage in Canadian newspapers
%==========================================================================
\landscape
\noindent\textbf{Supplementary Table S1.} Previous studies on climate change coverage in Canadian newspapers.
\vspace{0.5em}
{\footnotesize
\refstepcounter{table}
\label{tab:literature_review}
\begin{longtable}{p{5cm}p{6cm}p{4cm}p{3cm}p{3.5cm}}
\toprule
\textbf{Study} & \textbf{Newspaper(s)} & \textbf{Scope} & \textbf{Language} & \textbf{Time Period} \\
\midrule
\endfirsthead
\multicolumn{5}{c}{Supplementary Table~S1\ -- \textit{Continued from previous page}} \\
\toprule
\textbf{Study} & \textbf{Newspaper(s)} & \textbf{Scope} & \textbf{Language} & \textbf{Time Period} \\
\midrule
\endhead
\midrule
\multicolumn{5}{r}{\textit{Continued on next page}} \\
\endfoot
\bottomrule
\endlastfoot

\textcite{good_framing_2008} & The Globe and Mail & National & English & 2007 \\
& The National Post & National & English & \\
& The Toronto Star & National & English & \\
& The Calgary Herald & Regional & English & \\
& The Edmonton Journal & Regional & English & \\
& The Chronicle Herald & Regional & English & \\
& The Montreal Gazette & Regional & English & \\
& The Charlottetown Guardian & Regional & English & \\
& Halifax Daily News & Regional & English & \\
& The Ottawa Citizen & Regional & English & \\
& The Star Phoenix & Regional & English & \\
& The Telegram & Regional & English & \\
& The Times Colonist & Regional & English & \\
& The Vancouver Sun & Regional & English & \\
& The Winnipeg Sun & Regional & English & \\
& Yukon News & Regional & English & \\
\midrule
\textcite{rowe_comparing_2011} & The Globe and Mail & National & English & 2007--2008 \\
& The Toronto Star & National & English & \\
& The Calgary Herald & Regional & English & \\
& The Toronto Sun & Regional & English & \\
& The Montreal Gazette & Regional & English & \\
& The Ottawa Citizen & Regional & English & \\
\midrule
\textcite{young_representations_2011} & The Globe and Mail & National & English & 1988--2008 \\
& The National Post & National & English & \\
\midrule
\textcite{anne_difrancesco_seeing_2011} & The Globe and Mail & National & English & 2008 \\
& The National Post & National & English & \\
\midrule
\textcite{ahchong_anthropogenic_2012} & The Globe and Mail & National & English & 1988--2007 \\
& The Toronto Star & National & English & \\
\midrule
\textcite{young_comparing_2012} & The Globe and Mail & National & English & 2007--2008 \\
& The National Post & National & English & \\
& The Toronto Star & National & English & \\
& The Vancouver Sun & Regional & English & \\
& The Calgary Herald & Regional & English & \\
& The Telegram & Regional & English & \\
& Le Devoir & Regional & French & \\
& La Presse & Regional & French & \\
\midrule
\textcite{stoddart_canadian_2015} & The Globe and Mail & National & English & 1999--2010 \\
& The National Post & National & English & \\
\midrule
\textcite{ford_coverage_2015} & The Globe and Mail & National & English & 1993--2013 \\
& The Toronto Star & National & English & \\
\midrule
\textcite{stoddart_endangered_2016} & The Globe and Mail & National & English & 2006--2010 \\
& The National Post & National & English & \\
\midrule
\textcite{stoddart_media_2017} & The Globe and Mail & National & English & 1997--2010 \\
& The National Post & National & English & \\
\midrule
\textcite{barkemeyer_media_2017} & The Globe and Mail & National & English & 2008 \\
& The National Post & National & English & \\
& The Toronto Star & National & English & \\
& The Vancouver Sun & Regional & English & \\
& The Toronto Sun & Regional & English & \\
\midrule
\textcite{king_how_2019} & The Globe and Mail & National & English & 2005--2015 \\
& The National Post & National & English & \\
& The Toronto Star & National & English & \\
& The Calgary Herald & Regional & English & \\
& The Chronicle Herald & Regional & English & \\
& The Vancouver Sun & Regional & English & \\
& The Whitehorse Daily Star & Regional & English & \\
& Journal de Montréal & Regional & French & \\
\midrule
\textcite{pillod_reframing_2021} & The Globe and Mail & National & English & 2008--2020 \\
\midrule
\textcite{brandt_news_2021} & The Toronto Star & National & English & 2015--2018 \\
& The Calgary Herald & Regional & English & \\
& The Vancouver Sun & Regional & English & \\
& The Chronicle Herald & Regional & English & \\
& The Winnipeg Free Press & Regional & English & \\
& Times and Transcript & Regional & English & \\
\midrule
\textcite{hase_climate_2021} & The Globe and Mail & National & English & 2006--2018 \\
& The Toronto Star & National & English & \\
\midrule
\textcite{sachdeva_themes_2022} & n/a & n/a & English & 1986--2016 \\
\midrule
\textcite{stoddart_what_2023} & The Globe and Mail & National & English & 2015--2020 \\
& The National Post & National & English & \\
& The Telegraph & Regional & English & \\
& The Telegram & Regional & English & \\
& The Chronicle Herald & Regional & English & \\
& The Guardian & Regional & English & \\
\midrule
\textcite{chen_how_2023} & n/a & n/a & English & 2018--2021 \\
\bottomrule
\end{longtable}
}

\vspace{0.5em}
\noindent\footnotesize
Note: This table summarizes the newspaper sources, geographic scope, language, and time periods covered by previous studies on climate change media coverage in Canada. Most studies focus exclusively on English-language national newspapers, particularly The Globe and Mail and The National Post.
\endlandscape

%==========================================================================
% Supplementary Table S2 -- Newspapers included in the CCF database
%==========================================================================
\landscape
\noindent\textbf{Supplementary Table S2.} Newspapers included in the CCF database.
\vspace{0.5em}
{\footnotesize
\refstepcounter{table}
\label{tab:newspapers}
\begin{longtable}{llllr}
\toprule
\textbf{Province/Territory} & \textbf{Newspaper} & \textbf{Language} & \textbf{Avg. Weekday Circ. (2015)} & \textbf{Articles} \\
\midrule
\endfirsthead
\multicolumn{5}{c}{Supplementary Table~S2\ -- \textit{Continued from previous page}} \\
\toprule
\textbf{Province/Territory} & \textbf{Newspaper} & \textbf{Language} & \textbf{Avg. Weekday Circ. (2015)} & \textbf{Articles} \\
\midrule
\endhead
\midrule
\multicolumn{5}{r}{\textit{Continued on next page}} \\
\endfoot
\bottomrule
\endlastfoot

National & The Globe and Mail & English & 226,171 & 31,127 \\
 & The National Post & English & 79,157 & 20,829 \\
 & The Toronto Star & English & 192,084 & 48,616 \\
 & CBC.ca\textsuperscript{b} & English & -- & 2,941 \\
 & Radio-Canada.ca\textsuperscript{b} & French & -- & 3,054 \\
\midrule
Alberta & The Edmonton Journal & English & 59,244 & 18,763 \\
 & The Calgary Herald & English & 58,865 & 20,349 \\
\midrule
British Columbia & The Vancouver Sun & English & 78,901 & 19,068 \\
 & The Times Colonist & English & 47,556 & 12,111 \\
\midrule
Manitoba & The Winnipeg Free Press & English & 67,503 & 13,129 \\
\midrule
New Brunswick & L'Acadie Nouvelle\textsuperscript{c} & French & 18,102 & 5,143 \\
\midrule
Newfoundland \& Labrador & The Telegram & English & 15,380 & 6,208 \\
\midrule
Nova Scotia & The Chronicle Herald & English & 71,304 & 11,018 \\
\midrule
Ontario & The Toronto Sun & English & 85,775 & 3,333 \\
 & Le Droit\textsuperscript{c} & French & 28,616 & 4,728 \\
\midrule
Quebec & Le Devoir\textsuperscript{c} & French & 28,729 & 13,685 \\
 & Le Journal de Montréal & French & 175,220 & 5,790 \\
 & La Presse & French & 98,657 & 6,917 \\
 & La Presse+\textsuperscript{a} & French & -- & 11,636 \\
 & The Montreal Gazette & English & 53,344 & 9,799 \\
\midrule
Saskatchewan & The Star Phoenix & English & 31,081 & 8,018 \\
\midrule
Yukon & The Whitehorse Daily Star\textsuperscript{c} & English & 945 & 7,702 \\
\midrule
\textbf{Total} & & & & \textbf{283,964} \\
\end{longtable}
}

\vspace{0.5em}
\noindent\footnotesize\textsuperscript{a} In 2013, La Presse launched a free digital tablet edition called La Presse+ which offers slightly different content. The corpus considers articles published both on the website and in the digital tablet edition. Duplicates were discarded.\\
\textsuperscript{b} News websites of the Canadian public broadcaster, added to the corpus by the CCF Observatory; their coverage begins in 2023. No print circulation applies.\\
\textsuperscript{c} Coverage window bounded by the access period of the source aggregators: Le Droit until December 2023, The Whitehorse Daily Star until May 2024, Le Devoir and L'Acadie Nouvelle until March 2025. All other outlets are covered through July 2026 (corpus cut-off of release v\CCFversion{}).
\endlandscape

%==========================================================================
% Supplementary Table S3 -- Complete CCF annotation framework
%==========================================================================
\landscape
\noindent\textbf{Supplementary Table S3.} Complete CCF annotation framework.
\vspace{0.5em}
{\footnotesize
\phantomsection
\refstepcounter{table}
\label{tab:complete_framework}
\begin{longtable}{p{0.5cm}p{4.5cm}p{2.8cm}p{13.5cm}}
\toprule
\textbf{\#} & \textbf{Category} & \textbf{Code} & \textbf{Description} \\
\midrule
\endfirsthead
\multicolumn{4}{c}{Supplementary Table~S3\ -- \textit{Continued from previous page}} \\
\toprule
\textbf{\#} & \textbf{Category} & \textbf{Code} & \textbf{Description} \\
\midrule
\endhead
\midrule
\multicolumn{4}{r}{\textit{Continued on next page}} \\
\endfoot
\bottomrule
\endlastfoot

\multicolumn{4}{l}{\cellcolor{gray!10}\textbf{MAIN FRAMES}} \\
\midrule
\rowcolor{gray!8}
\multicolumn{4}{l}{\textit{Economic Frame}} \\
\midrule
1 & \textbf{Economic Frame\newline (Primary Category)} & \texttt{economic\_frame} & If the sentence refers to climate change as an economic issue, including its negative or positive impacts on the economy, the financial costs or benefits of climate action, or the carbon footprint of the economic or industrial sector. This label captures any framing of climate change in economic terms, such as effects on jobs, markets, industries, or public and private finances. \\
\cmidrule(lr){1-4}
2 & Negative impacts of climate change on the economy & \texttt{eco\_neg\_impact} & If the sentence refers to the negative effects of climate change on the economy, such as damage to infrastructure, loss of agricultural productivity, disruptions to trade, increased costs for businesses, or economic decline in sectors vulnerable to climate shifts. This includes any economic losses directly resulting from changing climate conditions. \\
3 & Positive impacts of climate change on the economy & \texttt{eco\_pos\_impact} & If the sentence refers to the positive effects of climate change on the economy, such as new agricultural opportunities, the opening of new trade routes, or the potential for economic growth in sectors adapting to climate shifts. This includes any economic gains directly resulting from changing climate conditions. \\
4 & Economic disadvantages of climate action & \texttt{eco\_cost} & If the sentence refers to the negative economic impacts caused by climate policies or climate action, such as costs associated with transitioning to green energy, investments in climate adaptation, debt incurred from financing climate action, or short-term job losses in fossil fuel industries. This category captures financial burdens and challenges linked to climate action or policy. \\
5 & Economic benefits of climate action & \texttt{eco\_benefit} & If the sentence refers to the positive economic impacts resulting from climate policies or climate action, such as job creation in renewable energy sectors, investments in green technologies, economic growth driven by sustainability initiatives, or cost savings from energy efficiency. This category captures financial gains and opportunities linked to climate action or policy. \\
6 & Carbon footprint of the economic sector & \texttt{eco\_footprint} & If the sentence refers to the environmental or carbon footprint of the economic or industrial sectors, including emissions and resource use related to manufacturing, transportation, energy production, mining, agriculture, and other commercial or industrial activities. This label covers the ecological impact generated by economic and industrial processes, not the effects of climate change on these sectors. \\
\midrule
\rowcolor{gray!8}
\multicolumn{4}{l}{\textit{Health Frame}} \\
\midrule
7 & \textbf{Health Frame\newline (Primary Category)} & \texttt{health\_frame} & If the sentence refers to climate change as a public health issue, including its negative or positive impacts on physical or mental health, the health co-benefits of climate action (e.g., cleaner air, active transportation), or the carbon footprint of the health sector. This label captures any framing of climate change in relation to health outcomes, healthcare systems, or health-related consequences and solutions. \\
\cmidrule(lr){1-4}
8 & Negative impacts of climate change on health & \texttt{health\_neg\_impact} & If the sentence refers to the negative effects of climate change on human health, such as increased risks of infectious diseases (e.g., malaria, dengue), respiratory and cardiovascular illnesses due to air pollution and heat exposure, heat-related illnesses (e.g., heatstroke), and mental health disorders (e.g., anxiety, depression) linked to environmental changes. \\
9 & \textit{Positive impacts of climate change on health}* & \texttt{health\_pos\_impact} & If the sentence refers to the potential positive effects of climate change on human health, such as reductions in cold-related illnesses and deaths, or increased opportunities for outdoor activity due to milder winters. \\
10 & Health co-benefits of climate action & \texttt{health\_cobenefit} & If the sentence refers to health benefits resulting from climate policies or actions, such as improved air quality, better nutrition, reduced respiratory or cardiovascular diseases, fewer premature deaths, and other positive health outcomes linked to climate mitigation and adaptation efforts. \\
11 & \textit{Carbon footprint of the health sector}* & \texttt{health\_footprint} & If the sentence refers to the environmental or carbon footprint of the health sector, including emissions and resource use associated with healthcare facilities, medical equipment, pharmaceuticals production, patient and staff travel, and waste generated by medical activities. This label covers the ecological impact of delivering health services, not the impacts of climate change on health outcomes. \\
\midrule
\rowcolor{gray!8}
\multicolumn{4}{l}{\textit{Security Frame}} \\
\midrule
12 & \textbf{Security Frame\newline (Primary Category)} & \texttt{security\_frame} & If the sentence refers to the effects of climate change on security, such as violent conflict driven by resource scarcity, an influx of climate refugees, disruption of military operations, or the carbon footprint of the defense sector. This category includes broader security challenges arising from climate change, including geopolitical tensions or the strategic responses to climate impacts. \\
\cmidrule(lr){1-4}
13 & Presence of climate refugees & \texttt{security\_refugees} & If the sentence refers to the displacement of populations or the presence of climate refugees as a result of a natural disaster, environmental degradation, or climate change impact. This category captures the movement of people fleeing climate-related disasters. \\
14 & Conflict & \texttt{security\_conflict} & If the sentence refers to the emergence or intensification of conflicts over natural resource exploitation due to climate change, such as disputes over water, land, or energy resources exacerbated by shifting climate patterns. This includes conflicts arising from scarcity or competition for increasingly stressed resources. \\
15 & \textit{Post-disaster military assistance}* & \texttt{security\_military} & If the sentence refers to the deployment of military forces as reinforcement after a natural disaster, such as wildfires, floods, or other extreme events. This includes military support in disaster response and recovery, like providing humanitarian aid or maintaining order during a crisis. \\
16 & \textit{Disruption of military operations}* & \texttt{security\_disruption} & If the sentence refers to the disruption of military operations due to a natural disaster affecting a military base or facilities. This includes cases where climate-induced events damage or disrupt military infrastructure, affecting defense readiness or response capabilities. \\
\midrule
\rowcolor{gray!8}
\multicolumn{4}{l}{\textit{Justice Frame}} \\
\midrule
17 & \textbf{Justice Frame\newline (Primary Category)} & \texttt{justice\_frame} & If the sentence frames climate change as a justice or moral issue, including references to who benefits or loses from climate action, differentiated responsibilities across countries or groups, unequal impacts of climate change, disparities in access to mitigation or adaptation measures, or concerns about intergenerational justice. This label captures ethical, equity-based, or fairness-oriented dimensions of climate discourse. \\
\cmidrule(lr){1-4}
18 & Winners and losers of climate action & \texttt{justice\_winners} & If the sentence refers to how climate policies or actions create winners and losers, by benefiting certain groups (e.g., vulnerable populations, Indigenous communities, developing countries) while disadvantaging others (e.g., fossil fuel workers, carbon-intensive industries, or developed nations facing steeper transitions). This label focuses on distributional outcomes of climate policy. \\
19 & Differentiated responsibility & \texttt{justice\_responsibility} & If the sentence refers to the unequal responsibility for causing climate change, especially in reference to the principle of common but differentiated responsibilities. This includes statements emphasizing that developed countries, multinational corporations, or high-emitting individuals have historically contributed more to global warming and should therefore bear a greater share of the mitigation burden. \\
20 & Unequal vulnerability to climate change & \texttt{justice\_vulnerability} & If the sentence highlights how different populations are unequally affected by the impacts of climate change, due to geographic, social, or economic vulnerabilities. This includes references to marginalized groups such as women, seniors, racialized populations, people in developing countries, or those living in high-risk zones (e.g., coastal or arid regions). This label focuses on impacts, not responsibility or capacity. \\
21 & Unequal access to climate action & \texttt{justice\_access} & If the sentence refers to unequal capacity to act on climate change, emphasizing that countries, companies, or individuals do not have equal financial, technical, or institutional means to take climate action. For example, small businesses may struggle to decarbonize compared to multinational firms, and low-income households may not afford sustainable alternatives. This label reflects inequality in ability to act, not impacts or blame. \\
22 & Intergenerational justice & \texttt{justice\_intergen} & If the sentence refers to intergenerational justice, meaning that current decisions and actions should not compromise the rights or well-being of future generations. Do not apply this label to sentences that merely mention age differences or generational groups (e.g., ``young people are protesting'') unless the statement clearly relates to moral responsibility across generations or the long-term consequences of climate inaction. \\
\midrule
\rowcolor{gray!8}
\multicolumn{4}{l}{\textit{Political Frame}} \\
\midrule
23 & \textbf{Political Frame\newline (Primary Category)} & \texttt{political\_frame} & If the sentence frames climate change as a political issue, including references to policy adoption or implementation, political debates or controversies, political positioning by parties or leaders on climate issues, and public opinion data related to climate change. This label captures the role of governance, decision-making, and political dynamics in climate discourse. \\
\cmidrule(lr){1-4}
24 & Policy action & \texttt{pol\_action} & If the sentence refers to the formal approval, adoption, or enactment of a policy, plan, or strategy by a government at the local, provincial, national, or international level. This includes the passing of new legislation, regulations, government programs, or official commitments related to climate or environmental issues (e.g., approval of a carbon pricing law, ratification of a renewable energy policy, adoption of a climate action framework). \\
25 & Political debate & \texttt{pol\_debate} & If the sentence refers to a political debate or discussion about climate change, particularly focusing on the differing views or disagreements on climate policies, the effectiveness of past actions, or the responsibility of different levels of government (e.g., federal vs. provincial). This includes disagreements about specific climate initiatives, such as proposed or enacted policies, laws, or funding allocations. \\
26 & Political positioning & \texttt{pol\_position} & If the sentence refers to the explicit stance or position of a politician or political party on climate change or climate-related policies. This includes statements where a political figure or party expresses their support or opposition to specific climate actions (e.g., carbon taxes, energy transitions, or climate mitigation strategies), or when they make electoral promises regarding climate policies. \\
27 & Public opinion data & \texttt{pol\_opinion} & If the sentence refers to data, surveys, polls, or studies measuring public opinion, attitudes, beliefs, or perceptions related to climate change, environmental policies, or climate action. This label captures quantitative or qualitative information on how the public views climate issues. \\
\midrule
\rowcolor{gray!8}
\multicolumn{4}{l}{\textit{Scientific Frame}} \\
\midrule
28 & \textbf{Scientific Frame\newline (Primary Category)} & \texttt{scientific\_frame} & If the sentence frames climate change as a scientific issue, including references to scientific research, discoveries, explanations, or assessments of climate change. This includes mentions of scientific consensus or controversy, as well as expressions of uncertainty or certainty in climate science. This label captures the role of scientific knowledge and expertise in understanding and communicating climate change. \\
\cmidrule(lr){1-4}
29 & Scientific debate & \texttt{sci\_debate} & If the sentence refers to a past or ongoing debate or disagreement within the scientific community about aspects of climate science or climate-related technologies. This includes contested interpretations of data, disputes about methodologies, or ethical debates surrounding proposed scientific or technological solutions such as geoengineering. It does not include general questioning of climate science or statements that challenge the validity or sufficiency of climate evidence. \\
30 & Popularisation or scientific discovery & \texttt{sci\_discovery} & If the sentence explains or describes climate science mostly grounded in the natural sciences---such as chemistry, physics, or biology---or presents scientific or technological discoveries related to climate change. This includes explanations of mechanisms (e.g., the greenhouse effect, climate feedbacks), climate modeling, or innovations such as carbon capture. It does not apply to statements that merely describe environmental conditions or trends without explaining an underlying scientific process. \\
31 & Questioning of climate science & \texttt{sci\_skepticism} & If the sentence questions the validity, accuracy, or sufficiency of climate science. This includes claims that climate projections are flawed, evidence is lacking, or conclusions are exaggerated. It only applies when science is being challenged or delegitimized, not when scientists acknowledge limitations, complexity, or areas of ongoing research. \\
32 & Defense of climate science & \texttt{sci\_defense} & If the sentence defends the robustness or credibility of climate science in response to skepticism or denial. This includes assertions that reinforce the scientific consensus, affirm the validity of climate models, or rebut claims that climate science is incorrect or misleading. It only applies when climate science is defended, not when it is neutrally presented. \\
\midrule
\rowcolor{gray!8}
\multicolumn{4}{l}{\textit{Environmental Frame}} \\
\midrule
33 & \textbf{Environmental Frame\newline (Primary Category)} & \texttt{environmental\_frame} & If the sentence frames climate change as an environmental issue, including references to the degradation or disappearance of natural habitats, loss of biodiversity, and negative impacts on fauna and flora. This label captures ecological consequences of climate change and emphasizes the relationship between climate disruption and the natural environment. \\
\cmidrule(lr){1-4}
34 & Loss of natural environments & \texttt{env\_habitat} & If the sentence refers to the long-term or irreversible disappearance or degradation of natural environments due to climate change. This includes situations where ecosystems or environmental features (e.g., glaciers, coral reefs) are predicted or expected to permanently disappear or significantly decline as a result of climate impacts, indicating that they will not recover. It does not apply to short-term, isolated, or recoverable events, such as one-time disasters (e.g., a single wildfire or a temporary drought) that might affect natural environments but do not necessarily imply permanent loss. \\
35 & Loss of fauna and flora & \texttt{env\_species} & If the sentence refers to the negative impacts of climate change on animal and plant species, including changes in habitat, shifts in migration or breeding patterns, increased mortality, loss of biodiversity, or threats to species survival. This label covers both temporary and ongoing effects on fauna and flora caused by climate change. \\
\midrule
\rowcolor{gray!8}
\multicolumn{4}{l}{\textit{Cultural Frame}} \\
\midrule
36 & \textbf{Cultural Frame\newline (Primary Category)} & \texttt{cultural\_frame} & If the sentence frames climate change as a cultural issue, including cultural depictions of climate change through art, literature, music, or film; challenges to hosting artistic or sports events due to climate impacts; the loss or disruption of Indigenous cultural practices; or the carbon footprint of the cultural and creative sectors. This label captures how climate change affects and is expressed through cultural life and identities. \\
\cmidrule(lr){1-4}
37 & Artistic representation & \texttt{cult\_art} & If the sentence refers to cultural representations of climate change in various forms such as art, literature, music, film, documentaries, theater, or other creative expressions. This label captures how climate change is portrayed, interpreted, or communicated through cultural and artistic mediums. \\
38 & Difficulty to host cultural or sports events & \texttt{cult\_event\_impact} & If the sentence refers to how climate change disrupts or challenges the organization of artistic, cultural, or sports events, such as cancelled festivals, concerts, or competitions due to extreme weather (e.g., storms, heatwaves, wildfires, or lack of snow for winter sports). This label captures climate-related obstacles to event planning or continuity. \\
39 & Loss of indigenous practices & \texttt{cult\_indigenous} & If the sentence refers to the erosion or disruption of traditional cultural practices among Indigenous communities due to climate change, such as changes in animal migration, plant cycles, or land access that affect food systems, ceremonies, or knowledge transmission. This label emphasizes climate-driven impacts on Indigenous cultural continuity. \\
40 & Carbon footprint of the cultural and sports sectors & \texttt{cult\_footprint} & If the sentence refers to the environmental or carbon footprint of the cultural, artistic, or sports sectors, including the emissions and resource use associated with activities such as touring, international travel, energy-intensive venues, unsustainable stage or set production, or event-related waste. This label covers the ecological impact of producing and hosting cultural or sports events, not the impacts of climate on those events. \\
\midrule
\multicolumn{4}{l}{\cellcolor{gray!10}\textbf{PRIMARY CATEGORIES}} \\
\midrule
\rowcolor{gray!8}
\multicolumn{4}{l}{\textit{Actors/Messengers}} \\
\midrule
41 & \textbf{Presence of Messengers\newline (Primary Category)} & \texttt{messenger} & If the sentence includes mentions or quotes of persons or organizations by journalists, whether through direct quotations or indirect references. This label captures the presence of sources, experts, officials, activists, or other actors cited or named within the text. A single individual or organization may be associated with multiple types of expertise or roles (e.g., a doctor would fall both under Health expert and Natural scientist). \\
\cmidrule(lr){1-4}
42 & Health expert & \texttt{msg\_health} & If the quoted or reported statement comes from a person or organization with medical or health-related recognized expertise. This includes individuals working in the fields of medicine, public health, healthcare administration, or biomedical research, whether in a professional, academic, or governmental capacity (e.g., professor of public health, doctor, nurse, family physician, public health officer, WHO representative, Red Cross medical coordinator, Doctors without Borders spokesperson, Canadian Medical Association representative, certified health consultants, minister of Health). \\
43 & Economic expert & \texttt{msg\_economic} & If the quoted or reported statement comes from a person or organization with economic or financial recognized expertise. This includes individuals working in the fields of economics, finance, business, or economic policy, whether in a professional, academic, private sector, or governmental capacity (e.g., economist, finance minister, central bank official, business executive, investment analyst, professor of economics, representative of a chamber of commerce, World Bank or IMF spokesperson, economic think tank expert, labour market analyst). \\
44 & Security expert & \texttt{msg\_security} & If the quoted or reported statement comes from a person or organization with recognized expertise in security or defence matters. This includes individuals with academic, professional, or operational experience in military, national security, intelligence, or strategic affairs (e.g., professor in defense or security studies, active or retired member of the armed forces, intelligence agency representative, conflict resolution specialist, national defence minister, security analyst, representative of a security-focused think tank or NGO). \\
45 & Legal expert & \texttt{msg\_legal} & If the quoted or reported statement comes from a person or organization with recognized expertise in law or legal affairs. This includes individuals working in legal professions, legal academics, the judiciary, or justice-related government roles (e.g., lawyer, prosecutor, law professor, minister of justice, attorney general, legal aid representative, bar association representative, legal consultants). \\
46 & Cultural or Sport expert & \texttt{msg\_cultural} & If the quoted or reported statement comes from a person with recognized expertise in sports, arts, or culture (e.g., Olympian, coach, professor of literature, museum director, filmmaker, writer, festival organizer, musician, artist, indigenous knowledge keeper). \\
47 & Natural scientist & \texttt{msg\_scientist} & If the quoted or reported statement comes from a person or organization with recognized expertise in the natural or hard sciences. This includes individuals working in fields such as climatology, environmental science, physics, chemistry, biology, earth sciences, engineering, or related disciplines, whether in academic, governmental, or research institutions (e.g., climate scientist, environmental researcher, oceanographer, physicist, ecologist, geologist, engineer, professor of atmospheric sciences, IPCC contributor, Environment Canada scientist, research institute spokesperson). \\
48 & Social scientist & \texttt{msg\_social} & If the quoted or reported statement comes from a person with expertise in social sciences, helping to understand how society perceives, responds to, and governs issues like climate change (e.g., sociologist, political scientist, media scholar, economist). \\
49 & Activist & \texttt{msg\_activist} & If the quoted or reported statement comes from a person or organization known primarily for advocacy or activism. This includes individuals or groups engaged in campaigning, protesting, or raising awareness on environmental, social, or climate-related issues, whether through grassroots organizing, NGOs, or international movements (e.g., climate activist, environmental campaigner, Indigenous land defender, Fridays for Future organizer, Greenpeace spokesperson, Extinction Rebellion member, citizen activist, youth climate leader, environmental justice advocate). \\
50 & Public official & \texttt{msg\_official} & If the quoted or reported statement comes from a politician, elected official, government representative, or public decision-maker at any level of government. This includes individuals holding political office or involved in policymaking, regulation, or governance, whether municipal, provincial, federal, or international (e.g., prime minister, environment minister, member of parliament, city councillor, premier, senator, government spokesperson, cabinet member, UN climate negotiator, provincial regulator, parliamentary committee chair). \\
\midrule
\rowcolor{gray!8}
\multicolumn{4}{l}{\textit{Events}} \\
\midrule
51 & \textbf{Presence of Events\newline (Primary Category)} & \texttt{event} & If the sentence refers to the occurrence or mention of one or more events. This includes extreme meteorological events (e.g., storms, floods, heatwaves), political or governmental events (e.g., elections, policy announcements, trials, judiciary decisions), meetings or conferences, publications, cultural events, protests, or other significant happenings. \\
\cmidrule(lr){1-4}
52 & Extreme meteorological event & \texttt{evt\_weather} & If the sentence refers to the occurrence of an extreme meteorological or climate-related event that causes or risks causing significant disruption, damage, or harm. This includes events such as heatwaves, floods, wildfires, hurricanes, droughts, storms, or other weather-related disasters, whether sudden or prolonged. \\
53 & Meeting & \texttt{evt\_meeting} & If the sentence refers to an official visit, conference, summit, forum, meeting, convention, or assembly that brings together people to discuss, negotiate, or present on specific topics at the local, national, or international levels (e.g., town hall meeting, premiers' conference, ministerial roundtable, state visit, UN conference, international forum on clean energy, G7 meeting, COP28, World Economic Forum). \\
54 & Publication & \texttt{evt\_publication} & If the sentence refers to the release or publication of a document such as a report, academic article, survey results, op-ed, policy brief, or white paper, often aimed at informing the public, guiding decision-making, or contributing to public discourse (e.g., IPCC report, Statistics Canada survey, university study on climate impacts, editorial on climate policy, NGO report on emissions, think tank publication). \\
55 & Election & \texttt{evt\_election} & If the sentence refers to the occurrence of an election at the local, provincial, national, or international level, including general elections, by-elections, leadership races, referenda, or electoral campaigns (e.g., federal election, municipal by-election, provincial leadership race, EU parliamentary elections). \\
56 & New policy & \texttt{evt\_policy} & If the sentence refers to the announcement or unveiling of a policy, plan, or strategy at the local, provincial, national, or international level, including new legislation, regulations, government programs, or official commitments, including those related to climate or environmental issues (e.g., carbon pricing plan, renewable energy strategy, climate action framework). \\
57 & Judiciary decision & \texttt{evt\_judiciary} & If the sentence refers to the occurrence of a trial, court ruling, legal proceeding, or judicial decision at any level of the judicial system, including constitutional challenges, environmental lawsuits, regulatory hearings, or Supreme Court decisions, including those related to climate or environmental issues (e.g., court ruling on carbon pricing, environmental class action lawsuit, constitutional challenge to climate legislation, tribunal decision on pipeline approval). \\
58 & Cultural or Sports event & \texttt{evt\_cultural} & If the sentence refers to the organization or hosting of a sports, artistic, or cultural event (e.g., Olympics games, national hockey tournament, local marathon, film screening, music concert, theatre performance, book launch, mural festival, photography exhibit). \\
59 & Protest & \texttt{evt\_protest} & If the sentence refers to the organization of a protest (e.g., climate strike, anti-pipeline protest, extinction rebellion action, anti-racism march, women's right rally, union-led protest for better wages). \\
\midrule
\rowcolor{gray!8}
\multicolumn{4}{l}{\textit{Solutions}} \\
\midrule
60 & \textbf{Presence of Solutions\newline (Primary Category)} & \texttt{solution} & If the sentence refers to solutions addressing climate change, including mitigation efforts to reduce greenhouse gas emissions or adaptation strategies to manage climate impacts. This includes direct mentions of specific policies, technologies, behaviors, or initiatives, as well as indirect references to approaches aimed at combating or coping with climate change. \\
\cmidrule(lr){1-4}
61 & Mitigation strategy & \texttt{sol\_mitigation} & If the sentence refers to mitigation solutions designed to reduce or prevent the emission of greenhouse gases and limit the extent of future climate change. These strategies aim to address the root causes of climate change by lowering carbon footprints, enhancing carbon sinks, or transitioning to low-carbon technologies (e.g., investing in renewable energy, improving energy efficiency, implementing carbon pricing, promoting public transit, reforestation, phasing out fossil fuels). \\
62 & Adaptation strategy & \texttt{sol\_adaptation} & If the sentence refers to adaptation solutions designed to reduce vulnerability to the negative effects of climate change by adjusting or preparing for those impacts. These strategies involve actions to manage the risks and enhance resilience, particularly by reducing exposure to climate-related hazards and minimizing harm to people, infrastructure, and ecosystems (e.g., providing cooling centres during heat waves, restoring wetlands to limit flooding, altering farming techniques to cope with changing weather patterns). \\
\midrule
\multicolumn{4}{l}{\cellcolor{gray!10}\textbf{EMOTIONAL TONE}} \\
63 & Emotion: positive & \texttt{tone\_positive} & If the tone is generally optimistic, reassuring, or enthusiastic, highlighting favorable aspects or expressing confidence in climate solutions or actions. \\
64 & Emotion: negative & \texttt{tone\_negative} & If the tone is generally alarming, critical, or pessimistic, especially if the sentence expresses concern, fear, or frustration about climate change or its impacts. \\
65 & Emotion: neutral & \texttt{tone\_neutral} & If the sentence is mainly informative or descriptive without an evaluative tone, presenting facts or balanced perspectives without strong emotional content. \\
\midrule
\multicolumn{4}{l}{\cellcolor{gray!10}\textbf{GEOGRAPHIC FOCUS}} \\
66 & Mention of Canada & \texttt{canada} & If the sentence refers to Canada in any context, domestic or international, including but not limited to climate change impacts on Canadian territory, as well as Canada's policies, actions, or roles. \\
\midrule
\multicolumn{4}{l}{\cellcolor{gray!10}\textbf{URGENCY/ALARMISM}} \\
67 & Urgency to act & \texttt{urgency} & If the sentence conveys a strong sense of urgency or alarmism, emphasizing the immediate need for action or highlighting critical risks and dangers related to climate change, often with a tone that signals warning or crisis. \\
\bottomrule
\end{longtable}
\vspace{0.5em}
\noindent\footnotesize
* Categories marked with an asterisk were included in the annotation protocol but had insufficient positive training examples. This table enumerates the 67 \emph{candidate} categories of the framework; the deployed database carries 65 of them. Two candidates (\texttt{health\_pos\_impact}, \texttt{health\_footprint}) had no positive examples in either language and were excluded entirely. Two deployed categories (\texttt{security\_military}, \texttt{security\_disruption}) could not be trained in English and are annotated on French sentences only (their English annotations are \texttt{NULL} by construction). The deployed model suite therefore counts $65 \times 2 - 2 = 128$ models. Because this table numbers the 67 candidates while \texttt{CCF\_reliability\_tiers} numbers the 65 deployed categories, the two numbering schemes differ from the Health frame onward: always match categories on the \texttt{code} column, never on the row number.
}
\endlandscape

%==========================================================================
% Supplementary Table S4 -- Complete model training performance metrics for all annotation categories.
%==========================================================================
\landscape
\noindent\textbf{Supplementary Table S4.} Complete model training performance metrics for all annotation categories.
\vspace{0.5em}

{\footnotesize
\input{../Results/Outputs/Tables/table_s4_training_metrics.tex}
}

\vspace{0.5em}
\noindent\footnotesize
\textit{Note.} F1 scores at the best epoch retained by the combined-criterion best-model selector ($0.98 \cdot F_1^{(1)} + 0.02 \cdot$ macro-$F_1$). EN columns report metrics for the English BERT-base-cased classifiers; FR columns report metrics for the French CamemBERT-base classifiers. A ``--'' in one language column means the category could not be trained in that language for lack of positive examples: this concerns exactly two English models (\texttt{security\_military}, \texttt{security\_disruption}), which is why the deployed suite counts $65 \times 2 - 2 = 128$ models; the two candidate categories excluded in both languages do not appear in this table (see the note to Supplementary Table~S3). The dataset distribution behind each row is reported in Supplementary Table~S5.

\endlandscape

%==========================================================================
% Supplementary Table S5 -- Training and validation dataset distribution across all annotation categories.
%==========================================================================
\landscape
\noindent\textbf{Supplementary Table S5.} Training and validation dataset distribution across all annotation categories.
\vspace{0.5em}

{\footnotesize
\input{../Results/Outputs/Tables/table_s5_training_distribution.tex}
}

\vspace{0.5em}
\noindent\footnotesize
\textit{Note.} Positive (\textbf{Pos}) and negative (\textbf{Neg}) sentence counts in the training and validation splits, by language. Categories with zero counts in every split could not be trained due to insufficient positive examples; the resulting per-category training metrics are reported in Supplementary Table~S4.

\endlandscape

%==========================================================================
% Supplementary Table S6 -- Named Entity Recognition performance
%==========================================================================
\noindent\textbf{Supplementary Table S6.} Named Entity Recognition performance.
\vspace{0.5em}
\begin{table}[h!]
\centering
\caption{Named Entity Recognition (\texttt{ner\_entities} variable) model performance metrics}
\label{tab:ner_performance}
\small
\begin{tabular}{llccc}
\toprule
\textbf{Language} & \textbf{Model} & \textbf{PER F1} & \textbf{ORG F1} & \textbf{LOC F1} \\
\midrule
English & BERT-base-NER & 0.961 & 0.811 & 0.925 \\
\midrule
\multirow{2}{*}{French} & spaCy fr\_core\_news\_lg & 0.880 & -- & -- \\
& CamemBERT-NER & -- & 0.824 & 0.929 \\
\bottomrule
\end{tabular}

\vspace{0.5em}
\noindent\footnotesize
Note: The hybrid approach for French combines spaCy for person entities (PER) with CamemBERT-NER for organization (ORG) and location (LOC) entities based on empirical evaluation on the dataset. Performance metrics are from the original model documentation.
\end{table}

%==========================================================================
% Supplementary Table S7 -- Detailed validation performance metrics on the 1{,}000-sentence gold standard.
%==========================================================================
\landscape
\noindent\textbf{Supplementary Table S7.} Detailed validation performance metrics on the 1{,}000-sentence gold standard.
\vspace{0.5em}

{\scriptsize
\input{../Results/Outputs/Tables/table_s7_test_metrics.tex}
}

\vspace{0.5em}
\noindent\footnotesize
\textit{Note.} F1 macro, F1 micro, F1 weighted, and positive-class support computed on the 1{,}000-sentence gold standard. EN, FR, and ALL columns report metrics on the English subset, the French subset, and their union. \textbf{OVERALL PERFORMANCE} in the last row reports the \emph{pooled} aggregates---all per-category predictions are pooled into a single confusion matrix before the scores are computed---and is therefore not the mean of the per-category rows. The unweighted mean of the 65 per-category macro F1 values (ALL column) is 0.773; this is the quantity compared like-for-like with the training-phase mean of 0.826 in the main text (Technical Validation). Support columns report the number of expert-confirmed positive examples; 23 of the 65 categories fall below the sampling target of 40 (see the main text on how to weigh low-support estimates). Inter-coder agreement statistics on the same 1{,}000 sentences are reported in Supplementary Tables~S9 and~S10.

\endlandscape

%==========================================================================
% Supplementary Table S8 -- Database-wide distribution of annotation categories across the \CCFsentApprox{}~million sentences of CCF\_processed\_data.
%==========================================================================
\landscape
\noindent\textbf{Supplementary Table S8.} Database-wide distribution of annotation categories across the \CCFsentApprox{}~million sentences of CCF\_processed\_data.
\vspace{0.5em}

{\footnotesize
\input{../Results/Outputs/Tables/table_s8_database_distribution.tex}
}

\vspace{0.5em}
\noindent\footnotesize
\textit{Note.} Counts and proportions are computed live from \texttt{CCF\_Database.CCF\_processed\_data} at table-build time; the ``ALL'' columns aggregate English and French. Proportions are expressed in percent of the total number of sentences in each language. The total sentence counts at the bottom of the table match the data-dictionary entry for \texttt{CCF\_processed\_data} in Supplementary Table~S12.

\endlandscape



%==========================================================================
% Supplementary Table S9 -- Per-category inter-coder reliability
%==========================================================================
\landscape
\noindent\textbf{Supplementary Table S9.} Per-category inter-coder reliability metrics on the 1{,}000-sentence gold standard, with the blind phase (sentences 601--1000) reported separately.
\vspace{0.5em}

{\footnotesize
\input{../Results/Outputs/Tables/table_s9_per_category_reliability.tex}
}

\vspace{0.5em}
\noindent\footnotesize
\textit{Note.} Cohen's $\kappa$, Krippendorff's $\alpha$, Gwet's AC1, and percent agreement are computed between the expert annotator (gold standard) and the independent second annotator. Only Cohen's $\kappa$ is reported in both forms---full-sample (all 1{,}000 doubly-coded sentences) and blind-phase; the other indices are reported for the blind phase only (sentences 601--1000, coded without any communication with the expert), which is the primary reliability benchmark used throughout the paper. Categories with ``--'' in the blind-phase columns had no positive examples in sentences 601--1000.

\endlandscape

%==========================================================================
% Supplementary Table S10 -- Reliability tiers
%==========================================================================
\landscape
\noindent\textbf{Supplementary Table S10.} Reliability tier assignment for each of the 65 annotation categories.
\vspace{0.5em}

{\footnotesize
\input{../Results/Outputs/Tables/table_s10_reliability_tiers.tex}
}

\vspace{0.5em}
\noindent\footnotesize
\textit{Note.} Each category is assigned to a reliability tier based on its validation $F_1$ macro and the blind-phase Cohen's $\kappa$. \textbf{Tier~A} (research-grade) requires $F_1 \geq 0.80$ \emph{and} $\kappa \geq 0.60$; \textbf{Tier~C} (exploratory) corresponds to $F_1 < 0.60$ \emph{or} $\kappa < 0.40$; \textbf{Tier~B} (use with caution) captures everything in between. The acronym \textbf{PIRD*} in the Flag column stands for \emph{prevalence-induced reliability deflation}: it identifies categories with $F_1 \geq 0.70$ but blind-phase Cohen's $\kappa < 0.40$, where the chance-corrected agreement is depressed by very low prevalence rather than by systematic disagreement between coders \autocite{feinstein_high_1990, byrt_bias_1993}. The two categories without an English model (\texttt{security\_military} and \texttt{security\_disruption}, annotated on French sentences only for lack of positive English training examples) are flagged ``no EN model'' and forced to Tier~C in every language column; the \texttt{exclusion\_reason} column of the deposited lookup carries the value \texttt{insufficient\_training\_data} for these two rows and is \texttt{NULL} everywhere else.

\endlandscape

%==========================================================================
% Supplementary Table S11 -- Per-model training provenance
%==========================================================================
\landscape
\noindent\textbf{Supplementary Table S11.} Per-model training provenance: best epoch, training phase, validation macro $F_1$, and positive-class support for every BERT/CamemBERT classifier.
\vspace{0.5em}

{\scriptsize
\input{../Results/Outputs/Tables/table_s11_training_hyperparameters.tex}
}

\vspace{0.5em}
\noindent\footnotesize
\textit{Note.} The English models are BERT-base-cased (case-sensitive: load their tokeniser with \texttt{do\_lower\_case=False}) and the French models are CamemBERT-base, both trained through the AugmentedSocialScientist framework. The four per-language sub-columns are: \textbf{Ep.} (best epoch retained by the combined-criterion best-model selector $0.98 \cdot F_1^{(1)} + 0.02 \cdot$ macro-$F_1$), \textbf{Ph.} (training phase --- $n$ for normal, $r$ for reinforced training, the latter firing when the positive-class $F_1$ falls below 0.70 during normal training), \textbf{$F_1$} (macro $F_1$ at the retained epoch), and \textbf{$n^{+}$} (number of positive validation examples; ``--'' denotes models that failed to train because of insufficient positive examples). The static configuration (optimiser AdamW with $\epsilon=10^{-8}$, learning rate $5\times 10^{-5}$ in normal training and $5\times 10^{-6}$ in the reinforced phase, batch size 32 normal / 64 reinforced, maximum sequence length 512 tokens, up to 20 epochs in each phase, random seed 42 for Python/NumPy/PyTorch RNGs, linear schedule with zero warm-up, an unweighted cross-entropy in normal training and, in the reinforced phase, a \texttt{WeightedRandomSampler} combined with a class-weighted cross-entropy (positive-class weight 3.5), training on an Apple Mac Studio M2~Ultra with 128~GB unified memory) is documented in full in the Methods section of the main manuscript.

\endlandscape

%==========================================================================
% Supplementary Table S12 -- Complete data dictionary
%==========================================================================
\landscape
\noindent\textbf{Supplementary Table S12.} Complete data dictionary of the enriched CCF Database: every column of every table, its type, and a one-line definition.
\vspace{0.5em}

{\footnotesize
\input{../Results/Outputs/Tables/table_s12_data_dictionary.tex}
}

\vspace{0.5em}
\noindent\footnotesize
\textit{Note.} All annotation column codes (\texttt{economic\_frame}, \texttt{eco\_neg\_impact}, etc.) and their definitions appear in Supplementary Table~S3.

\endlandscape

\printbibliography[title={Supplementary references}]

\end{document}
